\begin{pastexam}[label={px:gk3-107-ds-01}]{高考三級資訊處理 107 年 第一題}
（申論；25 分）

\begin{itemize}[leftmargin=2.2em]
\item[㈠] 請說明並比較二分搜尋（binary search）與一般二元搜尋樹（binary search tree）兩者在儲存鍵值並應用來
進行搜尋鍵值功能時，在'建置'與'搜尋'程序上作法與效能的差異（13 分）。
\item[㈡] 若有 n 個鍵值，以下列甲和乙兩種資料結構策略儲存：\\
策略甲：由小到大依序儲存在一陣列中\\
策略乙：以 AVL tree 架構儲存\\
請以 Big-O 觀念比較後續六種不同功能獨立運作時，這兩種策略何者效能較優或兩者效能相近：1.~尋找特定鍵值
k；2.~尋找排序為 j 的鍵值；3.~刪除特定鍵值 k；4.~刪除排序為 j 的鍵值；5.~插入新鍵值；6.~依序輸出所有鍵值。
（12 分）
\end{itemize}
\end{pastexam}

\begin{pastsoln}
答案：㈠ 二分搜尋先把鍵排序存在陣列（建置 $O(n \log n)$），搜尋穩定 $O(\log n)$；一般二元搜尋樹逐一插入
建置（平均 $O(n \log n)$、最壞 $O(n^2)$），搜尋平均 $O(\log n)$、最壞 $O(n)$，但插入刪除方便。 ㈡ 1 相近、
2 甲較優、3 乙較優、4 相近（AVL 加上子樹大小欄位時乙較優）、5 乙較優、6 相近。

\textbf{㈠ 建置。}
\begin{itemize}
\item \textbf{二分搜尋}：鍵存在一個連續的陣列，必須先由小到大排好。 用合併排序或堆積排序建置要
$O(n \log n)$；如果鍵已經排好，直接放進陣列只要 $O(n)$。 只需要 $n$ 格陣列，不需要指標。
\item \textbf{一般二元搜尋樹}：從空樹開始逐一插入（資結書虛擬碼~9.6），每次從根往下比到空的位置掛上去，
成本是當時的樹高。 鍵的順序隨機時，樹高期望 $O(\log n)$，建置期望 $O(n \log n)$（資結書第~9.3~節）；
但輸入已經排好（遞增或遞減）時，樹退化成一條鏈，第 $i$ 個要比 $i - 1$ 次，建置 $\Theta(n^2)$。 每個節點
多兩個指標欄位，空間較多。
\end{itemize}

\textbf{㈠ 搜尋。}
\begin{itemize}
\item \textbf{二分搜尋}：用索引直接算出中間位置 $\lfloor(\mathit{low} + \mathit{high})/2\rfloor$，每比一次範圍
減半，最多比 $\lfloor \lg n \rfloor + 1$ 次，\textbf{最壞也是} $O(\log n)$，因為每次都切在正中間。
\item \textbf{二元搜尋樹}：從根往下，比節點小往左、大往右（資結書虛擬碼~9.1），比較次數是目標的深度。
平均 $O(\log n)$，最壞（退化成鏈）$O(n)$；中間位置由樹的形狀決定，不一定在正中間。
\end{itemize}
\textbf{差異總結}：兩者搜尋的邏輯相同（都是「和中間的鍵比，決定往哪一半」），二分搜尋的「中間」是算出來
的、一定平衡；二元搜尋樹的「中間」是插入順序決定的。 二分搜尋的弱點在\textbf{更新}：陣列插入或刪除一個鍵
要搬動 $O(n)$ 個元素；二元搜尋樹只要改指標，成本是樹高。 所以資料固定不變、只查詢時用二分搜尋；常常插入
刪除時用二元搜尋樹（最好是平衡的，例如 AVL 樹）。

\textbf{㈡} AVL 樹的高度是 $O(\log n)$（資結書定理~10.1），所以搜尋、插入、刪除都是 $O(\log n)$。 下表的「乙」
先以一般的 AVL 樹（節點只存鍵、左右指標、高度或平衡因子）為準：
\begin{center}
\small
\begin{tabular}{lllc}
\toprule
功能 & 甲：排序陣列 & 乙：AVL 樹 & 較優 \\
\midrule
1. 尋找鍵 $k$ & $O(\log n)$ 二分搜尋 & $O(\log n)$ 從根往下 & 相近 \\
2. 尋找第 $j$ 小 & $O(1)$，就是 $A[j]$ & $O(n)$，中序走到第 $j$ 個 & 甲 \\
3. 刪除鍵 $k$ & $O(\log n)$ 找到，再搬 $O(n)$ 個 & $O(\log n)$ & 乙 \\
4. 刪除第 $j$ 小 & $O(1)$ 找到，再搬 $O(n)$ 個 & $O(n)$ 找到，$O(\log n)$ 刪 & 相近 \\
5. 插入新鍵 & $O(\log n)$ 找位置，再搬 $O(n)$ 個 & $O(\log n)$ & 乙 \\
6. 依序輸出全部 & $O(n)$ 從頭掃到尾 & $O(n)$ 中序走訪 & 相近 \\
\bottomrule
\end{tabular}
\end{center}
說明：
\begin{itemize}
\item 1：兩者都每一步把範圍縮小一半左右，$O(\log n)$。
\item 2：陣列可以用索引直接取 $A[j]$（從 $1$ 開始數時是 $A[j - 1]$）。 一般 AVL 樹不知道每棵子樹有幾個節點，
只能做中序走訪數到第 $j$ 個，$O(j + \log n)$，最壞 $O(n)$。
\item 3、5：陣列在 $O(\log n)$ 找到位置，但刪除要把後面的元素往前搬、插入要往後搬，平均搬 $n/2$ 個，$O(n)$。
AVL 樹刪除或插入後沿路徑往上調整平衡，最多 $O(\log n)$ 次旋轉，總共 $O(\log n)$。
\item 4：陣列找到第 $j$ 個是 $O(1)$，但刪除仍要搬 $O(n - j)$ 個；一般 AVL 樹要 $O(n)$ 才找到第 $j$ 個。 兩者
最壞都是 $O(n)$，相近。
\item 6：陣列本來就排好；AVL 樹做中序走訪，每個節點拜訪常數次，都是 $O(n)$。
\end{itemize}
\textbf{補充}：如果 AVL 樹的每個節點多存一個 $\mathit{size}$ 欄位（以它為根的子樹節點數，資結書第~9.4~節），
旋轉時 $O(1)$ 就能更新，第 $j$ 小可以用 select 在 $O(\log n)$ 找到。 這時 2 變成「甲 $O(1)$ 仍略優，兩者
都是對數以內」，4 變成乙 $O(\log n)$ 較優。 答題時要說明採用的是哪一種 AVL 樹。
\end{pastsoln}
